% Mushroom Livability - paper.
%
% Build:   cd paper && latexmk -pdf paper.tex      (or: pdflatex paper.tex, twice)
%
% Every number, table and figure is \input from paper/generated/, which is
% written by `python main.py pipeline` (or `python main.py paper`).
% Do not type results into this file by hand.
\documentclass[11pt,a4paper]{article}

\usepackage[utf8]{inputenc}
\usepackage[T1]{fontenc}
\usepackage{lmodern}
\usepackage[margin=2.5cm]{geometry}
\usepackage{amsmath}
\usepackage{booktabs}
\usepackage{xcolor}
\usepackage{tikz}
\usepackage[hidelinks]{hyperref}

\IfFileExists{generated/macros.tex}{\input{generated/macros.tex}}{%
  \GenericError{}{Missing generated/macros.tex}{Run `python main.py pipeline` (or `python main.py paper`) first.}{}}

\newcommand{\MLSscore}{Mushroom Livability Score}

\title{Where Should a Mushroom Lover Live?\\[4pt]
\large A transparent, open-data livability score for Paris, Amsterdam,\\ thirty other European cities and eight Canadian benchmarks}
\author{Pastel project}
\date{Data retrieved \MLSGeneratedDate}

\begin{document}
\maketitle

\ifMLSsynthetic
\begin{center}
\fcolorbox{red}{red!8}{\parbox{0.9\textwidth}{\centering\textbf{SYNTHETIC DEMO BUILD.}
The numbers in this document are pseudo-random placeholders produced by the demo mode.
They are not results. Re-run the pipeline without \texttt{--demo}.}}
\end{center}
\fi

\begin{abstract}
\noindent We ask where a person who loves finding wild mushrooms should live. We combine five ingredients
--- how many fungi are recorded, how many species, for how much of the year, how favourable the climate is
and how much woodland is nearby --- into a single \MLSscore{} between 0 and 100, computed identically for
\MLSNumLocations{} locations from open data only (iNaturalist, GBIF, ERA5 via Open-Meteo, WorldPop and
OpenStreetMap). Every score decomposes exactly into the contributions of its ingredients.
In the head-to-head that motivates the study, \MLSHeadlineWinner{} scores \MLSHeadlineWinnerScore{} and
\MLSHeadlineLoser{} scores \MLSHeadlineLoserScore{}. Across all \MLSNumScored{} scored locations the highest
score is \MLSLeaderScore{} (\MLSLeaderName{}, \MLSLeaderCountry{}); the best Canadian benchmark is
\MLSBestBenchmarkName{} (\MLSBestBenchmarkScore{}). The code, the cached API responses and this paper are
reproducible from a single command.
\end{abstract}

\section{Introduction}

``Where should I live?'' is usually answered with rents, salaries and commute times. For someone whose
weekends are spent looking for fungi the relevant questions are different: is there much to find near
home, is it varied, does the season last, and is there forest within reach? This paper turns those
questions into a small, auditable index.

The headline comparison is between Paris and Amsterdam. To give that comparison context, the same method
is applied to major cities across Germany and Central Europe, the Nordic countries, the United Kingdom and
Ireland, Southern Europe, France and the Benelux, and --- as a long-distance benchmark --- to eight
Canadian cities.

Two design principles guided the work. First, \emph{openness}: only public, openly licensed data are used,
and every request made to a data provider is stored with its URL and retrieval time. Second,
\emph{explainability}: the score is a weighted sum of a few interpretable quantities, so that any ranking
can be traced back to the numbers that produced it.

\section{Data}\label{sec:data}

A location is defined as the disc of radius \MLSRadiusKm{}~km (area \MLSAreaKm{}~km\textsuperscript{2})
around a city centre: roughly what can be reached on a half-day trip without a long journey.
All quantities below refer to that disc.

\paragraph{iNaturalist (primary).} Research-grade observations of the kingdom Fungi (iNaturalist taxon
47170, which includes lichens) are queried through the public API~\cite{inat}. We retrieve the total number
of fungal observations, the total number of research-grade observations of all taxa, the number of distinct
fungal taxa, the distribution of observations over the twelve months of the year, and the coordinates of
the \MLSSampleCap{} most recently added records (used only for the density maps in the web application).

\paragraph{GBIF (secondary).} The number of georeferenced occurrence records of kingdom Fungi without
geospatial issues is retrieved from the GBIF occurrence API~\cite{gbif}. GBIF adds museum specimens and
national recording schemes; it also contains iNaturalist's research-grade export, so the two sources
overlap and are not independent.

\paragraph{Climate.} Daily mean 2~m temperature and daily precipitation for \MLSClimatePeriod{} come from the
ERA5 reanalysis~\cite{era5} through the Open-Meteo historical weather API~\cite{openmeteo}, and are reduced
to twelve monthly normals (mean temperature, mean monthly precipitation total).

\paragraph{Population.} The number of residents inside the disc is obtained from the WorldPop
100~m gridded population surface for \MLSWorldPopYear{}~\cite{worldpop} using the WorldPop statistics service.

\paragraph{Habitat.} The share of the disc covered by woodland is estimated from
OpenStreetMap~\cite{osm} by point sampling: a regular lattice of points is tested for membership of
polygons tagged \texttt{landuse=forest} or \texttt{natural=wood}, and the share of points inside woodland
estimates the mapped woodland share.

\section{Method}\label{sec:method}

\subsection{Metrics}

Let $N_f$ be the number of research-grade fungal observations, $N_a$ the number of research-grade
observations of all taxa, $S$ the number of fungal taxa, $G$ the GBIF record count, $P$ the population and
$A$ the area of the disc. The metrics are
\begin{align*}
d_{\text{area}} &= N_f / A, & d_{\text{cap}} &= 10^5\, N_f / P, & d_{\text{share}} &= N_f / N_a, & d_{\text{gbif}} &= G / A .
\end{align*}
The per-capita rate keeps large cities from winning on headcount alone, and the fungal share corrects for
differences in how widely iNaturalist is used: a region with many users has many records of everything.

\paragraph{Seasonal reliability.} With $n_m$ the number of observations in calendar month $m$ and
$N_f=\sum_m n_m$, a month is \emph{reliable} when $n_m/N_f \ge 0.04$ and $n_m \ge 10$. The metric is the
number of reliable months $R\in\{0,\dots,12\}$. A perfectly even year would put $8.3\%$ in each month, so
the threshold asks for roughly half of an even share; the absolute floor prevents sparsely recorded
places from appearing reliable.

\paragraph{Climate suitability.} For each month with mean temperature $T_m$ (\textdegree C) and precipitation
$P_m$ (mm) we define
\[
 s_m = \tau(T_m)\cdot \min\!\left(1, \frac{P_m}{80}\right),\qquad
 \tau(T)=\begin{cases}
 0 & T\le 0 \ \text{or}\ T \ge 26\\
 T/8 & 0<T<8\\
 1 & 8\le T\le 18\\
 (26-T)/8 & 18<T<26
 \end{cases}
\]
and the climate suitability index is $C=\tfrac{1}{12}\sum_m s_m$. The window reflects the broad finding
that fruiting of most temperate macrofungi requires mild temperatures together with sufficient
moisture~\cite{boddy}; it is a deliberately simple heuristic, not a fitted model.

\paragraph{Habitat.} $F$ is the woodland share of the disc as described in Section~\ref{sec:data}.

\subsection{Normalisation and aggregation}

The count-like metrics ($d_{\text{area}}$, $d_{\text{cap}}$, $d_{\text{share}}$, $d_{\text{gbif}}$ and $S$) are heavy-tailed.
They are transformed with $\log(1+x)$ and then scaled linearly so that the weakest scored location receives
0 and the strongest 1. Bounded metrics are scaled against a fixed ceiling: $R/12$, $C/1$ and
$\min(1, F/0.6)$. Writing $\tilde{x}$ for a normalised metric, the score of a location is
\[
 \text{MLS} = 100 \sum_{c} w_c \sum_{m\in c} w_{c,m}\, \tilde{x}_{m},
\]
with component weights $w_c$ and within-component weights $w_{c,m}$ each summing to one.
Table~\ref{tab:model} lists the weights that were actually used in this run.

\begin{table}[ht]
\centering\footnotesize
\input{generated/model_table.tex}
\caption{Components, sub-metrics and effective weights (as a share of the total score).}
\label{tab:model}
\end{table}

\subsection{Missing data}

Three rules, applied before any arithmetic, keep scores comparable. (i)~A location lacking a required
metric ($d_{\text{area}}$, $S$, $R$ or $C$) is not scored. (ii)~An optional metric that is unavailable for
any scored location is dropped for all locations. (iii)~If a component loses all of its metrics it is
dropped and the remaining weights are rescaled proportionally. The outcome of these rules for the present
run is reported in Section~\ref{sec:coverage}.

\section{Results}

\subsection{Summary of findings}

\input{generated/findings.tex}

\subsection{Paris and Amsterdam}

Table~\ref{tab:h2h} compares the two headline cities ingredient by ingredient. Bold rows give the points
each component contributes to the total; the rows beneath show the underlying raw metrics.

\begin{table}[ht]
\centering\small
\input{generated/headtohead_table.tex}
\caption{Head-to-head comparison of the headline cities (higher-scoring city first).}
\label{tab:h2h}
\end{table}

\subsection{All locations}

Figure~\ref{fig:contrib} shows how the leading scores are composed and Table~\ref{tab:ranking} gives the
complete ranking.

\begin{figure}[ht]
\centering
\input{generated/contributions_figure.tex}
\caption{Composition of the \MLSscore{} for the highest-ranked locations and the headline cities.
Segment lengths are points contributed by each component.}
\label{fig:contrib}
\end{figure}

\begin{table}[p]
\centering\footnotesize\setlength{\tabcolsep}{4.5pt}
\input{generated/ranking_table.tex}
\caption{Full ranking. Headline cities in bold, Canadian benchmarks in italics. Component columns are
points contributed to the score.}
\label{tab:ranking}
\end{table}

\subsection{Coverage of this run}\label{sec:coverage}

\input{generated/coverage.tex}

\section{Discussion}

% interpretation.tex - the only hand-written results prose in the paper.
%
% RULES FOR WHOEVER EDITS THIS FILE (human or agent) - see AGENTS.md:
%   1. Write it only AFTER a full, real (non --demo) pipeline run.
%   2. Every claim must be checkable against data/processed/results.json.
%      Do not type numbers here: refer to tables/figures or use the \MLS...
%      macros from generated/macros.tex.
%   3. No causal claims the data cannot support. "Is recorded more often"
%      is fine; "has more mushrooms" is not.
%   4. Treat every location and region as a peer: no city or country is a
%      reference point. 2-4 short paragraphs: (a) what separates the two
%      highest-ranked locations, (b) what the leaders have in common,
%      (c) how the regions compare (regions in alphabetical order).
%   5. Then delete the placeholder paragraph below.
%
% Rewritten 2026-11-21 for the location-neutral presentation, on the results of the
% 2026-11-19 real run (results.json: synthetic=false, partial=false; scores unchanged).
% Every statement below can be checked against Tables~\ref{tab:h2h}, \ref{tab:regions} and \ref{tab:ranking}.

\paragraph{The two highest-ranked locations.} \MLSLeaderName{} and \MLSRunnerUpName{} win on different
ingredients (Table~\ref{tab:h2h}). \MLSLeaderName{} is ahead on seasonal reliability, with every month of the
year counted as reliable, and on species richness. \MLSRunnerUpName{} is ahead on observation density. Within
that component the iNaturalist rates of the two cities are close; the difference lies in GBIF records per
km\textsuperscript{2}, where \MLSRunnerUpName{} is far ahead. GBIF depth reflects which recording schemes
publish their data, so this part of the gap describes recording and publication effort as much as fungal life.
Climate suitability and woodland contribute similar points to both.

\paragraph{The leaders.} The five highest-ranked locations come from five different regions
(Figure~\ref{fig:contrib}, Table~\ref{tab:regions}), so no single part of the study area dominates the top of
the ranking. What the leaders share is strength on the record-based ingredients --- observation density,
species richness and seasonal reliability --- rather than on climate: \MLSLeaderName{} and \MLSRunnerUpName{}
both receive fewer climate points than most of the field, while the climate index favours mild Atlantic
locations, several of which rank in the lower half. Because the leaders are large, well-observed cities, this
pattern is consistent with the observer-effort caveat below, and the score should be read as ``recorded fungal
life'', not as abundance.

\paragraph{Regions.} Differences within regions are as large as differences between them
(Table~\ref{tab:ranking}). In Canada, in France and the Benelux, and in Germany and Central Europe, the
locations span the ranking from the top five to the bottom three, so the region a city lies in says little
about its score. The best location of every region except Southern Europe ranks in the top five overall
(Table~\ref{tab:regions}). The Southern European locations all rank in the middle or lower part of the table;
their median points for observation density and species richness are the lowest of any region, while their
median climate and seasonal-reliability points are close to the median of all locations. With citizen-science data,
lower recording density can reflect fewer recorders as much as fewer fungi.

\paragraph{Caveats specific to this run.} Two data limitations shape these results (Section~\ref{sec:coverage}).
First, the WorldPop statistics service was unavailable, so the per-capita observation rate was dropped for all
locations; the density component in this run rests on area rates, fungal share and GBIF depth only. Second, the
woodland share counts only woodland that the Overpass API exposes as an area, which in practice means named
multipolygons. It is therefore a lower bound whose bias varies with local mapping practice: several plainly
forested locations, including Trento and Victoria, receive no habitat points. Cross-region comparisons of the
habitat ingredient should not be relied on, but it carries the smallest weight in the score.


\subsection{Limitations}

\begin{itemize}
\item \textbf{Observer effort.} Citizen-science records measure where people look as much as where fungi
grow~\cite{isaac}. The fungal-share and per-capita metrics reduce, but do not remove, this bias. A high
score means ``many fungi are recorded here'', which is evidence of, not identical to, ``many fungi grow here''.
\item \textbf{Relative scaling.} Count-like metrics are scaled between the weakest and strongest location in
the comparison set, so adding or removing locations changes scores (though rarely the order of two
given cities).
\item \textbf{Source overlap.} GBIF includes iNaturalist's research-grade export; the GBIF metric is
therefore partly redundant with the iNaturalist metrics and carries a small weight.
\item \textbf{Heuristic climate model.} The temperature and moisture windows are simple expert-style
heuristics applied to monthly normals; they ignore humidity, soil, altitude within the disc and year-to-year
variability.
\item \textbf{Habitat completeness.} OpenStreetMap woodland mapping varies in completeness between
countries, and the lattice estimate has a sampling error of several percentage points.
\item \textbf{Weights.} The component weights are a transparent judgement, not an estimate. Because
contributions are reported separately, readers can re-weight them.
\item \textbf{Scope.} Lichens are included in the iNaturalist Fungi taxon. Edibility, access rights and
foraging regulations are outside the scope of this study.
\end{itemize}

\section{Reproducibility}

\begin{sloppypar}
The project is a plain Python package with two runtime dependencies. \texttt{python main.py pipeline}
retrieves all data (responses are cached together with their URL and retrieval time), computes the scores
and regenerates every table, figure and number in this paper; \texttt{python main.py} then serves an
interactive web application with density maps and per-location breakdowns. Methodological parameters live
in \texttt{src/settings.py} and the scoring model in \texttt{src/scoring/components.py}.
\end{sloppypar}

\begin{thebibliography}{9}
\bibitem{inat} iNaturalist. Research-grade observations, accessed through the iNaturalist API v1.
\url{https://www.inaturalist.org}. Retrieved \MLSGeneratedDate.
\bibitem{gbif} GBIF.org. GBIF Occurrence Search API, kingdom Fungi. \url{https://www.gbif.org}.
Retrieved \MLSGeneratedDate.
\bibitem{era5} Hersbach, H., et al. (2020). The ERA5 global reanalysis. \emph{Quarterly Journal of the
Royal Meteorological Society}, 146(730), 1999--2049.
\bibitem{openmeteo} Zippenfenig, P. (2023). Open-Meteo.com Weather API. Zenodo.
\url{https://open-meteo.com}.
\bibitem{worldpop} WorldPop (\url{https://www.worldpop.org}), School of Geography and Environmental Science,
University of Southampton. Global high-resolution population denominators, 2000--2020.
\bibitem{osm} OpenStreetMap contributors. \url{https://www.openstreetmap.org/copyright}. Data available
under the Open Database Licence.
\bibitem{boddy} Boddy, L., et al. (2014). Climate variation effects on fungal fruiting.
\emph{Fungal Ecology}, 10, 20--33.
\bibitem{isaac} Isaac, N. J. B., and Pocock, M. J. O. (2015). Bias and information in biological records.
\emph{Biological Journal of the Linnean Society}, 115(3), 522--531.
\end{thebibliography}

\end{document}
